\section{Introduction}

Language-model systems can answer from parameters, retrieve external evidence, cache reusable facts, or update themselves temporarily or durably. Existing work usually studies retrieval selection \citep{asai2024selfrag,yan2024crag}, editing and memory mechanisms \citep{wang2024wise,yu2023melo}, or broad surveys of knowledge editing \citep{zhang2024keStudy} in isolation. We instead study a controller problem: when should one policy choose among parametric answering, retrieval, memory reads, memory writes, temporary adaptation, and guarded durable consolidation under a single objective?

The answer is not a universal-win story. Our frozen results show one clear positive regime and several clear negative ones. On PopQA recurring-memory \citep{mallen2022trust}, delayed utility is genuinely useful: the full router matches retrieval quality while cutting retrieval calls substantially, and removing delayed utility collapses performance. On public FreshQA-style changing-answer evaluation derived from FreshLLMs weekly releases \citep{vu2024freshllms}, retrieval dominates across all four models, and suppressing delayed utility is more reliable than softly calibrating it. On MQuAKE \citep{zhong2023mquake}, retrieval again dominates, and the gap persists on direct edits, single-hop probes, and multi-hop consequences. On UniEdit \citep{chen2025uniedit}, nearly every method is already at ceiling, so the benchmark serves mainly as a sanity check rather than a separator.

Those results support a sharper claim than the project originally targeted: delayed utility helps in stable recurring-memory settings, but hurts in public changing-answer and edit-propagation settings; in those regimes, retrieval-dominant behavior wins, and removing delayed utility is more reliable than softly calibrating it.

The public benchmark trust story is also intentionally modest. Our manual FreshQA audit found 24 clean cases, 76 suspect cases, and no direct leakage-risk cases. That rules out a simple leakage explanation for the retrieval-dominant public result, but it does not fully clear the benchmark. Appendix~A includes a small deterministic controlled-update diagnostic that is consistent with the same boundary, but the headline evidence in the paper comes from the frozen benchmark surface.

The paper makes four contributions:
\begin{itemize}[leftmargin=1.5em]
\item a unified six-action routing formulation over parameters, retrieval, reusable memory, temporary adaptation, and guarded durable consolidation;
\item a clean positive result on stable recurring-memory PopQA, where delayed utility preserves quality while reducing repeated retrieval;
\item a cross-benchmark boundary result showing that retrieval dominates public changing-answer and propagation-heavy regimes, with soft calibration offering only limited repair relative to outright suppression of delayed utility;
\item a frozen evidence package spanning FreshQA public main and sequence runs, MQuAKE, UniEdit, a manual FreshQA audit, and benchmark reliability notes.
\end{itemize}
